\chapter{Methodology}
\label{sec:methodology}

\section{Methodology for Multi-Suction Cup Gripper Optimization}

This section describes the optimization framework developed for a multi-suction
cup gripper consisting of three suction cups arranged in a triangular
configuration. The framework leverages the multi-objective optimization library
\textit{pymoo}~\cite{pymoo} and delegates objective function evaluation to an
external service via an API call, enabling a modular and decoupled architecture.

\subsection{Parameterization of the Gripper}

Each suction cup is defined by its planar position, while the height and
orientation remain fixed. Thus, the decision variables are:

\begin{equation}
    \mathbf{p} = [(x_1, y_1),\; (x_2, y_2),\; (x_3, y_3)]
\end{equation}
\nomenclature{$\mathbf{p}$}{Decision variable vector representing the gripper configuration}
\nomenclature{$(x_i, y_i)$}{Planar position of the $i$-th suction cup}

Each individual $\mathbf{p}$ forms a triangular base of the gripper.

\paragraph{Assumption}
A triangular configuration of suction cups is assumed to yield a stable grasp~\cite{madan2022realrobotchallenge2021}.
This is motivated by the fact that three spatially distributed contact points
form a statically stable support, enabling effective load distribution and
resistance to external disturbances.

\subsection{Decision Variables and Bounds}

Each decision variable is defined with a type (continuous or integer) and
associated bounds:

\begin{equation}
    x_i \in [\ell_i,\; u_i], \quad i = 1, \dots, n
\end{equation}
\nomenclature{$\ell_i$}{Lower bound of the $i$-th decision variable}
\nomenclature{$u_i$}{Upper bound of the $i$-th decision variable}
\nomenclature{$n$}{Number of decision variables}

The variable definitions are provided by the user as part of the optimization
request, allowing the framework to handle mixed-variable problems (i.e.,
combinations of continuous and integer variables).

\subsection{Optimization Framework}

The optimization is formulated as a multi-objective minimization problem and
solved using the \textit{pymoo} library. The user configures the optimization
through a request specifying:

\begin{itemize}
    \item Population size $N$
    \item Maximum number of generations $G_{\max}$
    \item Choice of optimization algorithm
    \item Number of objectives $n_{\text{obj}}$
    \item Decision variable definitions (type and bounds)
    \item Evaluation endpoint URL
\end{itemize}

\subsubsection{Supported Algorithms}

The framework supports the following multi-objective evolutionary algorithms:

\begin{table}[htbp]
\centering
\caption{Supported multi-objective optimization algorithms.}
\label{tab:algorithms}
\begin{tabular}{ll}
\toprule
\textbf{Identifier} & \textbf{Algorithm} \\
\midrule
NSGA2   & Non-dominated Sorting Genetic Algorithm II~\cite{nsga2} \\
NSGA3   & Non-dominated Sorting Genetic Algorithm III~\cite{nsga3} \\
MOEAD   & Multi-Objective Evolutionary Algorithm based on Decomposition~\cite{moead} \\
CMOPSO  & Competitive Multi-Objective Particle Swarm Optimization \\
SMSEMOA & S-Metric Selection Evolutionary Multi-Objective Algorithm~\cite{smsemoa} \\
\bottomrule
\end{tabular}
\end{table}

For algorithms requiring reference directions (NSGA-III, MOEA/D), uniform
reference directions are generated using Das and Dennis's
method~\cite{das_dennis}:

\begin{equation}
    \Lambda = \texttt{get\_reference\_directions}(\text{``uniform''},\; n_{\text{obj}},\; n_{\text{partitions}})
\end{equation}
\nomenclature{$\Lambda$}{Set of reference directions for decomposition-based algorithms}

For mixed-variable problems, specialized sampling and mating operators are
employed to handle heterogeneous variable types.

\subsection{Initialization}

An initial population of $N$ candidate solutions is generated via mixed-variable
sampling:

\begin{equation}
    \mathcal{P}^{(0)} = \{\mathbf{p}_1, \mathbf{p}_2, \dots, \mathbf{p}_N\}
\end{equation}
\nomenclature{$\mathcal{P}^{(0)}$}{Initial population of candidate solutions}
\nomenclature{$N$}{Population size}

\subsection{Objective Function Evaluation via API}

The evaluation of objective function values is decoupled from the optimization
loop and delegated to an external evaluation service. For each candidate
solution $\mathbf{p}$, the optimizer issues an API call to the evaluation
endpoint, which executes the grasp evaluation subprocess (described in
Section~\ref{sec:grasp_eval}) and returns the objective values.

\begin{equation}
    \mathbf{f}(\mathbf{p}) = \texttt{API\_call}(\mathbf{p}) \rightarrow [f_1(\mathbf{p}),\; f_2(\mathbf{p})]
\end{equation}

This architecture enables:
\begin{itemize}
    \item Separation of concerns between optimization logic and domain-specific evaluation
    \item Scalability through independent deployment of the evaluation service
    \item Reusability of the optimization framework for different problem domains
\end{itemize}

\subsubsection{Grasp Evaluation Subprocess}
\label{sec:grasp_eval}

The evaluation service receives a candidate gripper configuration $\mathbf{p}$
and performs the following steps:

\begin{enumerate}

\item \textbf{Grasp Generation:}
For each point in the workpiece point cloud, a candidate suction grasp is
generated by aligning the tool center point (TCP) of the gripper with the
surface point.

\item \textbf{Collision Filtering:}
The gripper is modeled as a rigid 3D body with fixed height, forming a
triangular prism whose base is defined by $\mathbf{p}$. Each candidate grasp is
checked for collisions between the gripper body and the workpiece. Colliding
grasps are discarded.

\item \textbf{Clustering of Valid Grasps:}
Let $\mathcal{G}_{\text{valid}}$ denote the set of collision-free grasps. If
$|\mathcal{G}_{\text{valid}}| \geq K_{\min}$, clustering is performed based on
spatial proximity and orientation similarity to reduce redundancy.
\nomenclature{$\mathcal{G}_{\text{valid}}$}{Set of valid (collision-free) grasps}
\nomenclature{$K_{\min}$}{Minimum number of valid grasps required for clustering}

\item \textbf{Objective Computation:}
The service computes and returns the objective values:

\begin{equation}
    \mathbf{f}(\mathbf{p}) = [-D(\mathbf{p}),\; -A(\mathbf{p})]
\end{equation}

where:
\begin{itemize}
    \item $A(\mathbf{p})$ is the triangle area (maximized for stability):
    \begin{equation}
        A(\mathbf{p}) = \frac{1}{2} \left| (x_2 - x_1)(y_3 - y_1) - (x_3 - x_1)(y_2 - y_1) \right|
    \end{equation}
    \item $D(\mathbf{p})$ is the diversity metric (maximized for robustness),
    computed from the spatial distribution and orientation spread of clustered grasps.
\end{itemize}

\nomenclature{$A(\mathbf{p})$}{Area of the triangle formed by the three suction cups}
\nomenclature{$D(\mathbf{p})$}{Diversity metric of the grasp distribution}
\nomenclature{$\mathbf{f}(\mathbf{p})$}{Multi-objective function vector}

Since \textit{pymoo} minimizes by convention, maximization objectives are
negated.

\end{enumerate}

\subsection{Optimization Loop}

The evolutionary optimization proceeds as follows:

\begin{equation}
    \mathcal{P}^{(t+1)} = \mathcal{A}\left(\mathcal{P}^{(t)},\; \mathbf{f}\right), \quad t = 0, 1, \dots
\end{equation}
\nomenclature{$\mathcal{A}$}{Selected multi-objective evolutionary algorithm}
\nomenclature{$t$}{Generation index}

where $\mathcal{A}$ denotes the selected algorithm performing selection,
crossover, mutation, and survival operations.

\subsection{Termination Criteria}

The optimization terminates when any of the following conditions is satisfied:

\begin{table}[htbp]
\centering
\caption{Termination criteria.}
\label{tab:termination}
\begin{tabular}{lll}
\toprule
\textbf{Criterion} & \textbf{Symbol} & \textbf{Description} \\
\midrule
Design space tolerance  & $x_{\text{tol}}$ & Change in decision variables below threshold \\
Objective space tolerance & $f_{\text{tol}}$ & Relative improvement in objectives below threshold \\
Maximum generations     & $G_{\max}$       & Maximum number of generations reached \\
Maximum evaluations     & $E_{\max}$       & Maximum number of function evaluations reached \\
\bottomrule
\end{tabular}
\end{table}
\nomenclature{$x_{\text{tol}}$}{Design space convergence tolerance}
\nomenclature{$f_{\text{tol}}$}{Objective space convergence tolerance}
\nomenclature{$G_{\max}$}{Maximum number of generations}
\nomenclature{$E_{\max}$}{Maximum number of function evaluations}

Tolerances are evaluated over a sliding window of $w$ generations.

\subsection{Performance Monitoring}

The hypervolume indicator~\cite{hypervolume} is tracked at each generation via a
callback mechanism:

\begin{equation}
    \text{HV}^{(t)} = \text{HV}(\mathcal{F}^{(t)},\; \mathbf{r})
\end{equation}
\nomenclature{$\text{HV}^{(t)}$}{Hypervolume indicator at generation $t$}
\nomenclature{$\mathbf{r}$}{Reference point for hypervolume computation}
\nomenclature{$\mathcal{F}^{(t)}$}{Objective values of the population at generation $t$}

where $\mathbf{r} \in \mathbb{R}^{n_{\text{obj}}}$ is the reference point and
$\mathcal{F}^{(t)}$ denotes the set of objective vectors at generation $t$. A
monotonically increasing hypervolume indicates convergence toward the Pareto
front.

\subsection{Optional Hyperparameter Optimization}

The framework optionally supports automated tuning of algorithm hyperparameters
(e.g., crossover probability, mutation rates) using a meta-optimization
approach. A \texttt{HyperparameterProblem} is constructed around the selected
algorithm and evaluated over multiple random seeds. The mean hypervolume across
runs serves as the meta-objective:

\begin{equation}
    \min_{\boldsymbol{\theta}} \left\{ -\frac{1}{|\mathcal{S}|} \sum_{s \in \mathcal{S}} \text{HV}(\mathcal{F}_s,\; \mathbf{r}) \right\}
\end{equation}
\nomenclature{$\boldsymbol{\theta}$}{Algorithm hyperparameters}
\nomenclature{$\mathcal{S}$}{Set of random seeds for meta-optimization}

The best hyperparameters are then applied to the main optimization run.

\subsection{Output}

The optimization produces a set of Pareto-optimal solutions:

\begin{equation}
    \mathcal{P}^* = \left\{ \mathbf{p} \;\middle|\; \nexists\; \mathbf{p}' \in \mathcal{P} : \mathbf{f}(\mathbf{p}') \preceq \mathbf{f}(\mathbf{p}) \right\}
\end{equation}
\nomenclature{$\mathcal{P}^*$}{Set of Pareto-optimal solutions}

The final output includes:
\begin{itemize}
    \item The Pareto-optimal set $\mathcal{P}^*$ with corresponding objective values
    \item The hypervolume progression $\{\text{HV}^{(0)}, \text{HV}^{(1)}, \dots, \text{HV}^{(T)}\}$
    \item Total optimization wall-clock time
\end{itemize}

These solutions represent optimal trade-offs between grasp diversity and
triangle area (stability).